% Source: Weinberg GR, physical PDF pp. 296--298 (printed pp. 274--276).
% Coverage: Section 10.6, Scattering and Absorption of Gravitational
%           Radiation; equations (10.6.1)--(10.6.15).
% Figures/tables/footnotes: reference markers 5--11; no figures or tables.
% Status: source-reviewed and compile-clean.
% Uncertainties: none.

\subsection{Scattering and Absorption of Gravitational Radiation}\label{sec:10.6}

Consider a plane \emph{gravitational} wave with polarization \(e_{\mu\nu}\)
and wave vector \(k^\mu\), impinging on a target at the origin. At great
distances from the target, the gravitational wave will in general consist of
the plane wave and an outgoing scattered wave,%
\hyperref[ch10-ref-5]{\textsuperscript{5}}
\begin{equation}
  h_{\mu\nu}(\mathbf{x},t)
  \xrightarrow{r\to\infty}
  \left[
    e_{\mu\nu}e^{\ii\mathbf{k}\cdot\mathbf{x}}
    +f_{\mu\nu}(\widehat{\mathbf{x}})\frac{e^{\ii\omega r}}{r}
  \right]e^{-\ii\omega t},
  \tag{10.6.1}\label{eq:10.6.1}
\end{equation}
where \(r=\lvert\mathbf{x}\rvert\),
\(\widehat{\mathbf{x}}=\mathbf{x}/r\),
\(\omega=\lvert\mathbf{k}\rvert\), and \(f_{\mu\nu}\) is a
\emph{scattering amplitude}, which may depend on
\(\widehat{\mathbf{x}}\) and \(\omega\), but not on \(r\) or \(t\).

In order to analyze the energy balance between the gravitational wave and the
target, it is necessary to decompose the wave \eqref{eq:10.6.1} into incoming
and outgoing parts. The plane-wave part has the Legendre expansion%
\hyperref[ch10-ref-6]{\textsuperscript{6}}
\[
  e^{\ii\mathbf{k}\cdot\mathbf{x}}
  =
  \sum_{l=0}^{\infty}
  (2l+1)P_l(\widehat{\mathbf{k}}\cdot\widehat{\mathbf{x}})
  \ii^l j_l(\omega r),
\]
where \(P_l\) is the usual Legendre polynomial and \(j_l\) is the spherical
Bessel function%
\hyperref[ch10-ref-7]{\textsuperscript{7}}
of order \(l\). Asymptotically, we have%
\hyperref[ch10-ref-8]{\textsuperscript{8}}
\[
  \ii^lj_l(\omega r)
  \longrightarrow
  \frac{1}{2\ii\omega r}
  \left[e^{\ii\omega r}-(-1)^le^{-\ii\omega r}\right],
\]
so the sums over \(l\) become simply the Legendre expansions of delta
functions:%
\hyperref[ch10-ref-9]{\textsuperscript{9}}
\[
  \sum_l(2l+1)P_l(\mu)=2\delta(1-\mu),
\]
\[
  \sum_l(2l+1)(-1)^lP_l(\mu)=2\delta(1+\mu).
\]
The plane wave may therefore be asymptotically decomposed into outgoing and
ingoing waves
\[
  e^{\ii\mathbf{k}\cdot\mathbf{x}}
  \xrightarrow{r\to\infty}
  \frac{e^{\ii\omega r}}{\ii\omega r}
  \delta(1-\widehat{\mathbf{k}}\cdot\widehat{\mathbf{x}})
  -
  \frac{e^{-\ii\omega r}}{\ii\omega r}
  \delta(1+\widehat{\mathbf{k}}\cdot\widehat{\mathbf{x}}),
\]
and the gravitational wave \eqref{eq:10.6.1} has the corresponding
decomposition
\begin{equation}
  h_{\mu\nu}
  \xrightarrow{r\to\infty}
  \left[
    e^{\mathrm{out}}_{\mu\nu}e^{\ii\omega r}
    +e^{\mathrm{in}}_{\mu\nu}e^{-\ii\omega r}
  \right]e^{-\ii\omega t}
  +\text{c.c.},
  \tag{10.6.2}\label{eq:10.6.2}
\end{equation}
where
\begin{equation}
  e^{\mathrm{out}}_{\mu\nu}(\widehat{\mathbf{x}})
  =
  \frac{1}{\ii\omega r}
  \left[
    e_{\mu\nu}
    \delta(1-\widehat{\mathbf{k}}\cdot\widehat{\mathbf{x}})
    +\ii\omega f_{\mu\nu}(\widehat{\mathbf{x}})
  \right],
  \tag{10.6.3}\label{eq:10.6.3}
\end{equation}
\begin{equation}
  e^{\mathrm{in}}_{\mu\nu}(\widehat{\mathbf{x}})
  =
  -\frac{1}{\ii\omega r}e_{\mu\nu}
  \delta(1+\widehat{\mathbf{k}}\cdot\widehat{\mathbf{x}}).
  \tag{10.6.4}\label{eq:10.6.4}
\end{equation}

The total power carried \emph{out} of a large sphere of radius \(r\) by the
outgoing wave part of \eqref{eq:10.6.2} may be calculated, following the
same reasoning as in Section~\ref{sec:10.4}, as
\begin{equation}
  P_{\mathrm{out}}
  =
  \int d\Omega\,
  \left\langle t_{\mathrm{out}}^{0i}\right\rangle
  \widehat x_i r^2,
  \tag{10.6.5}\label{eq:10.6.5}
\end{equation}
where \(\langle t_{\mathrm{out}}^{0i}\rangle\) is the mean energy flux, obtained
by using \(e^{\mathrm{out}}_{\mu\nu}\) in place of \(e_{\mu\nu}\) in
Eq.~\eqref{eq:10.3.5}, and averaging over spacetime dimensions large
compared with \(1/\omega\) and small compared with \(r\). Inspection of
Eq.~\eqref{eq:10.6.3} shows that \(P_{\mathrm{out}}\) will consist of three
terms,
\begin{equation}
  P_{\mathrm{out}}
  =
  P_{\mathrm{scat}}+P_{\mathrm{int}}+P_{\mathrm{plane}},
  \tag{10.6.6}\label{eq:10.6.6}
\end{equation}
which arise, respectively, from \(f_{\mu\nu}\) alone, from the interference
between \(f_{\mu\nu}\) and \(e_{\mu\nu}\), and from \(e_{\mu\nu}\) alone.
The first term, which represents the total power scattered away from the
incident direction, is calculated by using \eqref{eq:10.3.5} in
\eqref{eq:10.6.5}, with \(f_{\mu\nu}/r\) in place of \(e_{\mu\nu}\):
\begin{equation}
  P_{\mathrm{scat}}
  =
  \frac{\omega^2}{16\pi G}
  \int d\Omega
  \left[
    f^{\lambda\nu*}(\widehat{\mathbf{x}})
    f_{\lambda\nu}(\widehat{\mathbf{x}})
    -\frac12
    \left|f^\lambda{}_\lambda(\widehat{\mathbf{x}})\right|^2
  \right].
  \tag{10.6.7}\label{eq:10.6.7}
\end{equation}
The interference term is similarly calculated as
\[
  P_{\mathrm{int}}
  =
  \frac{\omega^2}{8\pi G}
  \operatorname{Re}\!\left\{
    -\frac{1}{\ii\omega}
    \int d\Omega\,
    \delta(1-\widehat{\mathbf{k}}\cdot\widehat{\mathbf{x}})
    \left[
      e^{\lambda\nu*}f_{\lambda\nu}(\widehat{\mathbf{x}})
      -\frac12(e^\lambda{}_\lambda)^*
      f^\nu{}_\nu(\widehat{\mathbf{x}})
    \right]
  \right\},
\]
or, integrating over the delta function,
\begin{equation}
  P_{\mathrm{int}}
  =
  -\frac{\omega}{4G}
  \operatorname{Im}\!\left\{
    e^{\lambda\nu*}f_{\lambda\nu}(\widehat{\mathbf{k}})
    -\frac12(e^\lambda{}_\lambda)^*
    f^\nu{}_\nu(\widehat{\mathbf{k}})
  \right\}.
  \tag{10.6.8}\label{eq:10.6.8}
\end{equation}
The last term in \eqref{eq:10.6.6}, which represents the power carried out
of the sphere by the plane wave, is formally infinite for \(r\to\infty\).
However, a plane wave carries as much power into any volume as it carries out
of it, so the power brought into a sphere of radius \(r\) by the ingoing wave
\eqref{eq:10.6.4} is also equal to this term:
\begin{equation}
  P_{\mathrm{in}}=P_{\mathrm{plane}}.
  \tag{10.6.9}\label{eq:10.6.9}
\end{equation}
Thus \(P_{\mathrm{plane}}\) cancels out of the equation of energy conservation,
which gives the power absorbed by the target as
\begin{equation}
  P_{\mathrm{abs}}
  =
  P_{\mathrm{in}}-P_{\mathrm{out}}
  =
  -P_{\mathrm{scat}}-P_{\mathrm{int}}.
  \tag{10.6.10}\label{eq:10.6.10}
\end{equation}
The energy flux in the incident wave is given by Eq.~\eqref{eq:10.3.5} as
\begin{equation}
  \Phi
  \equiv
  \left\langle t^{0i}\right\rangle\widehat k_i
  =
  \frac{\omega^2}{16\pi G}
  \left[
    e^{\lambda\nu*}e_{\lambda\nu}
    -\frac12\left|e^\nu{}_\nu\right|^2
  \right].
  \tag{10.6.11}\label{eq:10.6.11}
\end{equation}
Thus the effective cross-section for elastic scattering of the gravitational
wave is
\begin{equation}
  \begin{split}
  \sigma_{\mathrm{scat}}
  &\equiv\frac{P_{\mathrm{scat}}}{\Phi}\\
  &=
  \frac{
    \displaystyle\int d\Omega
    \left[
      f^{\lambda\nu*}(\widehat{\mathbf{x}})
      f_{\lambda\nu}(\widehat{\mathbf{x}})
      -\frac12
      \left|f^\lambda{}_\lambda(\widehat{\mathbf{x}})\right|^2
    \right]
  }{
    \displaystyle
    e^{\lambda\nu*}e_{\lambda\nu}
    -\frac12\left|e^\nu{}_\nu\right|^2
  }.
  \end{split}
  \tag{10.6.12}\label{eq:10.6.12}
\end{equation}
This must be distinguished from the \emph{total} cross-section for scattering
or absorption of the wave
\begin{equation}
  \sigma_{\mathrm{tot}}
  \equiv
  \frac{P_{\mathrm{scat}}+P_{\mathrm{abs}}}{\Phi}.
  \tag{10.6.13}\label{eq:10.6.13}
\end{equation}
According to \eqref{eq:10.6.10}, the total cross-section can be expressed in
terms of the interference between the incident and scattered waves,
\begin{equation}
  \sigma_{\mathrm{tot}}=-\frac{P_{\mathrm{int}}}{\Phi},
  \tag{10.6.14}\label{eq:10.6.14}
\end{equation}
or, using \eqref{eq:10.6.8} and \eqref{eq:10.6.11},
\begin{equation}
  \sigma_{\mathrm{tot}}
  =
  \frac{4\pi}{\omega}
  \frac{
    \operatorname{Im}\!\left\{
      e^{\lambda\nu*}f_{\lambda\nu}(\widehat{\mathbf{k}})
      -\frac12(e^\lambda{}_\lambda)^*
      f^\nu{}_\nu(\widehat{\mathbf{k}})
    \right\}
  }{
    e^{\lambda\nu*}e_{\lambda\nu}
    -\frac12\left|e^\lambda{}_\lambda\right|^2
  }.
  \tag{10.6.15}\label{eq:10.6.15}
\end{equation}
This result, that the total cross-section is \(4\pi/\omega\) times the
imaginary part of a forward scattering amplitude, was first derived in
classical electrodynamics,%
\hyperref[ch10-ref-10]{\textsuperscript{10}}
and is therefore known as the \emph{optical theorem}. Here and in
electrodynamics it is a consequence of the conservation of energy, whereas
in quantum mechanics there is a similar theorem based on the conservation of
probability.%
\hyperref[ch10-ref-11]{\textsuperscript{11}}

Since the incident wave is weak, the scattering amplitude \(f_{\mu\nu}\) is
a linear combination of the components of the incident polarization tensor
\(e_{\rho\sigma}\). It follows that the cross-sections
\eqref{eq:10.6.12} and \eqref{eq:10.6.15} are independent of the
normalization of \(e_{\mu\nu}\), though they may depend on \(k\) and on the
\emph{form} of the polarization tensor. The aim of gravitational scattering
theory is to calculate \(f_{\mu\nu}\); following this, the various
cross-sections can be determined from \eqref{eq:10.6.12} and
\eqref{eq:10.6.15}.
